\documentclass[12pt]{article}
\usepackage{fullpage}
\usepackage[T1]{fontenc}
\usepackage[utf8]{inputenc}
\usepackage{lmodern}
\usepackage{microtype}
\usepackage{amsmath,amssymb}
\usepackage{amsthm}
\usepackage{graphicx}
\usepackage{booktabs}
\usepackage{hyperref}
\usepackage{url}
\usepackage{color}
\usepackage{xcolor}
\usepackage{enumitem}
\usepackage{amsfonts}
\usepackage{algorithm}
\usepackage{algpseudocode}

\DeclareMathOperator{\vecop}{vec}
\DeclareMathOperator{\diag}{diag}
\DeclareMathOperator{\range}{range}
\DeclareMathOperator{\nullsp}{null}
\newcommand{\bR}{\mathbb{R}}
\newcommand{\tZ}{\tilde Z}
\newcommand{\tK}{\tilde K}

\newtheorem{theorem}{Theorem}
\newtheorem{lemma}{Lemma}
\newtheorem{proposition}{Proposition}
\newtheorem{remark}{Remark}

\title{Preconditioned CG for the RKHS-Constrained CP Mode Subproblem}
\author{}
\date{}

\begin{document}
\maketitle

\section{Setup and notation}

We solve the mode-$k$ subproblem of an RKHS-constrained CP decomposition with
missing data. With $N=\prod_i n_i$, $n=n_k$, $M=\prod_{i\ne k}n_i$ and $q$ the
number of observed entries, we are given
\begin{itemize}[nosep]
\item a symmetric positive semidefinite kernel matrix $K\in\bR^{n\times n}$;
\item the Khatri--Rao product $Z\in\bR^{M\times r}$;
\item the mode-$k$ unfolding $T\in\bR^{n\times M}$ of the (zero-filled) tensor;
\item the selection matrix $S\in\bR^{N\times q}$ picking the $q$ observed entries
from $\vecop(T)$;
\item the MTTKRP $B=TZ\in\bR^{n\times r}$.
\end{itemize}
The unknown is $W\in\bR^{n\times r}$ (so the RKHS factor is $A_k=KW$), and the
normal equations are
\begin{equation}\label{eq:system}
  H\,\vecop(W)=g,\qquad
  H:=(Z\otimes K)^{T}SS^{T}(Z\otimes K)+\lambda\,(I_r\otimes K),\qquad
  g:=(I_r\otimes K)\,\vecop(B),
\end{equation}
a linear system of size $nr\times nr$. We assume $d<n,r<q\ll N$ and $n\ll M$.
Throughout, $\Pi:=SS^{T}\in\bR^{N\times N}$ is the diagonal $0/1$ matrix that zeroes
all unobserved coordinates.

We index the observed entries by $\ell=1,\dots,q$. The $\ell$-th observed entry of
the mode-$k$ unfolding sits in row $i_\ell\in\{1,\dots,n\}$ and column
$j_\ell\in\{1,\dots,M\}$ of $T$; its value is $t_\ell:=T[i_\ell,j_\ell]$. We write
\[
  \tZ\in\bR^{q\times r},\quad \tZ[\ell,:]:=Z[j_\ell,:],
  \qquad
  \tK\in\bR^{q\times n},\quad \tK[\ell,:]:=K[i_\ell,:],
\]
for the row subsets of $Z$ and $K$ indexed by the observations; these are exactly
the objects denoted $\tilde Z,\tilde K$ in the problem statement. The vectorization
convention is column stacking, so the $(i,j)$ entry of an $n\times M$ matrix occupies
position $i+n(j-1)$ of its $\vecop$.

\section{Structural identities}

The efficiency of the method rests on three exact identities. We use repeatedly the
standard Kronecker--vectorization identity
\begin{equation}\label{eq:vec}
  \vecop(AXB)=(B^{T}\otimes A)\,\vecop(X),
\end{equation}
valid for conformable real matrices $A,X,B$ (Henderson--Searle); for $A=K$, $X=W$,
$B=Z^{T}$ it gives $(Z\otimes K)\vecop(W)=\vecop\!\big(KWZ^{T}\big)$.

\begin{lemma}[Model evaluation at observed entries]\label{lem:model}
For any $W\in\bR^{n\times r}$,
\[
  S^{T}(Z\otimes K)\,\vecop(W)=m\in\bR^{q},
  \qquad
  m_\ell=\big(K[i_\ell,:]\,W\big)\cdot Z[j_\ell,:]
        =\sum_{b=1}^{r}\big(KW\big)[i_\ell,b]\,\tZ[\ell,b].
\]
\end{lemma}
\begin{proof}
By \eqref{eq:vec}, $(Z\otimes K)\vecop(W)=\vecop(KWZ^{T})$, whose $(i,j)$ entry is
$(KWZ^{T})[i,j]=\big(K[i,:]W\big)\cdot Z[j,:]$. The matrix $S^{T}$ extracts the
coordinates $i+n(j-1)$ for $(i,j)=(i_\ell,j_\ell)$, $\ell=1,\dots,q$, which are
exactly the stated values. The last equality is the definition of the matrix product
$(KW)$ and of $\tZ$.
\end{proof}

\begin{lemma}[Adjoint scatter]\label{lem:adjoint}
For any $y\in\bR^{q}$,
\[
  (Z\otimes K)^{T}S\,y=\vecop\!\big(K\,R_y\,Z\big)\in\bR^{nr},
  \qquad
  R_y:=\sum_{\ell=1}^{q}y_\ell\,e_{i_\ell}e_{j_\ell}^{T}\in\bR^{n\times M},
\]
where $e_i$ denotes a standard basis vector. Moreover the $n\times r$ matrix
$R_yZ=\sum_{\ell}y_\ell\,e_{i_\ell}\,\tZ[\ell,:]$ can be formed without referencing
$Z$ outside its observed rows.
\end{lemma}
\begin{proof}
$Sy=\vecop(R_y)$ because $Sy$ is the vector with value $y_\ell$ at the coordinate
$i_\ell+n(j_\ell-1)$ and $0$ elsewhere, which is precisely $\vecop$ of
$R_y=\sum_\ell y_\ell e_{i_\ell}e_{j_\ell}^{T}$. Then
$(Z\otimes K)^{T}=(Z^{T}\otimes K^{T})=(Z^{T}\otimes K)$ since $K=K^{T}$, and by
\eqref{eq:vec} with $A=K$, $X=R_y$, $B=Z$,
\[
  (Z^{T}\otimes K)\vecop(R_y)=\vecop(K\,R_y\,Z).
\]
Finally $R_yZ=\sum_\ell y_\ell\,e_{i_\ell}\,(e_{j_\ell}^{T}Z)
=\sum_\ell y_\ell\,e_{i_\ell}\,Z[j_\ell,:]=\sum_\ell y_\ell\,e_{i_\ell}\,\tZ[\ell,:]$,
which uses only $\tZ$.
\end{proof}

\begin{proposition}[Matrix--vector product without $M$ or $N$]\label{prop:matvec}
For any $W\in\bR^{n\times r}$, writing $m\in\bR^q$ as in Lemma~\ref{lem:model},
\begin{equation}\label{eq:matvec}
  H\,\vecop(W)=\vecop\!\Big(K\,P\,+\,\lambda\,KW\Big),
  \qquad
  P:=\sum_{\ell=1}^{q}m_\ell\,e_{i_\ell}\,\tZ[\ell,:]\ \in\bR^{n\times r}.
\end{equation}
This product depends on the data only through $K\in\bR^{n\times n}$,
$\tZ\in\bR^{q\times r}$ and the index list $\{i_\ell\}_{\ell=1}^q$.
\end{proposition}
\begin{proof}
By definition $H=(Z\otimes K)^{T}SS^{T}(Z\otimes K)+\lambda(I_r\otimes K)$. Apply
Lemma~\ref{lem:model} to obtain $y:=S^{T}(Z\otimes K)\vecop(W)=m$. Apply
Lemma~\ref{lem:adjoint} with this $y=m$ to obtain
$(Z\otimes K)^{T}Sm=\vecop(K R_m Z)$ with $R_mZ=\sum_\ell m_\ell e_{i_\ell}\tZ[\ell,:]=P$.
For the regularizer, \eqref{eq:vec} with $A=K$, $X=W$, $B=I_r$ gives
$(I_r\otimes K)\vecop(W)=\vecop(KW)$. Summing the two contributions yields
\eqref{eq:matvec}. Every quantity appearing ($KW$, $m$, $P$, $KP$, $KW$) is computed
from $K$, $\tZ$, and $\{i_\ell\}$ only; none touches the $M$ unobserved columns or
the $N$-dimensional ambient space.
\end{proof}

\begin{proposition}[Right-hand side and MTTKRP without $M$ or $N$]\label{prop:rhs}
The MTTKRP satisfies
$B=TZ=\sum_{\ell=1}^{q}t_\ell\,e_{i_\ell}\,\tZ[\ell,:]\in\bR^{n\times r}$, and the
right-hand side of \eqref{eq:system} is $g=\vecop(KB)$.
\end{proposition}
\begin{proof}
Since $T$ is zero outside the observed entries, $T=\sum_\ell t_\ell e_{i_\ell}e_{j_\ell}^{T}$,
so $TZ=\sum_\ell t_\ell e_{i_\ell}(e_{j_\ell}^{T}Z)=\sum_\ell t_\ell e_{i_\ell}\tZ[\ell,:]$,
using only $\tZ$ and the values $t_\ell$. By \eqref{eq:vec} with $A=K$, $X=B$, $B=I_r$,
$g=(I_r\otimes K)\vecop(B)=\vecop(KB)$.
\end{proof}

\section{Symmetry, solvability, and the singular-kernel case}

\begin{lemma}[$H$ is symmetric PSD]\label{lem:psd}
$H$ in \eqref{eq:system} is symmetric and positive semidefinite. It is positive
definite if and only if $K$ is nonsingular.
\end{lemma}
\begin{proof}
Write $G:=(Z\otimes K)^{T}SS^{T}(Z\otimes K)=\big(S^{T}(Z\otimes K)\big)^{T}\big(S^{T}(Z\otimes K)\big)$,
a Gram matrix, hence symmetric PSD. Since $K=K^{T}\succeq0$ and $I_r\otimes K$ is
symmetric PSD, $H=G+\lambda(I_r\otimes K)\succeq0$ for $\lambda>0$ and is symmetric.

Let $x=\vecop(W)$. If $Hx=0$ then $x^{T}Hx=\|S^{T}(Z\otimes K)x\|^{2}+\lambda\,x^{T}(I_r\otimes K)x=0$,
forcing $x^{T}(I_r\otimes K)x=0$. Because $I_r\otimes K\succeq0$, this gives
$(I_r\otimes K)x=0$, i.e.\ $\vecop(KW)=0$, i.e.\ every column of $W$ lies in
$\nullsp(K)$. Conversely any such $W$ also satisfies $G x=0$ since
$(Z\otimes K)x=\vecop(KWZ^{T})=0$. Hence
\begin{equation}\label{eq:null}
  \nullsp(H)=I_r\otimes\nullsp(K)
  =\{\vecop(W): \text{every column of }W\in\nullsp(K)\},
\end{equation}
which is $\{0\}$ iff $\nullsp(K)=\{0\}$, i.e.\ iff $K$ is nonsingular.
\end{proof}

\begin{lemma}[Consistency and minimum-norm solution]\label{lem:consistent}
The system \eqref{eq:system} is always consistent. When $K$ is singular it has a
unique solution in $\range(H)=\range(I_r\otimes K)$, and conjugate gradients started
at $0$ converges to exactly that solution.
\end{lemma}
\begin{proof}
$K$ is symmetric, so $\bR^{n}=\range(K)\oplus\nullsp(K)$ orthogonally, hence
$\bR^{nr}=\range(I_r\otimes K)\oplus(I_r\otimes\nullsp(K))$ orthogonally. By
\eqref{eq:null}, $\range(H)=\nullsp(H)^{\perp}=\big(I_r\otimes\nullsp(K)\big)^{\perp}
=\range(I_r\otimes K)$. The right-hand side $g=\vecop(KB)$ lies in
$\range(I_r\otimes K)=\range(H)$, so the system is consistent. The CG iterates
started from $x_0=0$ remain in the Krylov subspace
$\mathrm{span}\{g,Hg,H^{2}g,\dots\}\subseteq\range(H)$, on which $H$ acts as a
symmetric positive-definite operator (its only zero eigenvalues are confined to
$\nullsp(H)\perp\range(H)$). Therefore CG converges to the unique solution lying in
$\range(H)$, which is the minimum-Euclidean-norm solution. (When $K$ is nonsingular,
Lemma~\ref{lem:psd} makes $H$ SPD and the solution is unique.)
\end{proof}

\section{Preconditioner}

We precondition with the Kronecker-structured operator
\begin{equation}\label{eq:precond}
  P_{\mathrm{c}}:=\big(\tZ^{T}\tZ+\lambda I_r\big)\otimes K
  \;=\;C\otimes K,\qquad C:=\tZ^{T}\tZ+\lambda I_r\in\bR^{r\times r}.
\end{equation}
Here $C\succ0$ (it is $\lambda I_r$ plus a Gram matrix), and $C$ uses only the
observed rows $\tZ$, costing $O(qr^{2})$ to form. The matrix $C\otimes K$ is the
exact regularizer block $\lambda(I_r\otimes K)$ \emph{plus} the Gram term obtained by
replacing the kernel cross-products with the identity, i.e.\ a Kronecker
factorization of the dominant structure of $H$; equivalently it is the natural
``observed-data'' Kronecker approximation. It is symmetric PSD and shares the
nullspace of $H$:
\[
  \nullsp(P_{\mathrm{c}})=\nullsp(C\otimes K)=I_r\otimes\nullsp(K)=\nullsp(H),
\]
because $C$ is nonsingular. Thus $P_{\mathrm{c}}$ and $H$ are simultaneously SPD on
the common invariant range $\range(I_r\otimes K)$ where CG operates
(Lemma~\ref{lem:consistent}), so $P_{\mathrm{c}}$ is a legitimate
(symmetric positive-definite, on $\range H$) preconditioner.

\begin{proposition}[Preconditioner solve without $M$ or $N$]\label{prop:precsolve}
Let $K=V\Sigma V^{T}$ be an eigendecomposition with $\Sigma=\diag(\sigma_1,\dots,\sigma_n)$,
$\sigma_i\ge0$, and set $\sigma_i^{+}=1/\sigma_i$ if $\sigma_i>\tau$ and $0$
otherwise ($\tau$ a small tolerance, $K^{+}=V\diag(\sigma^{+})V^{T}$ the
Moore--Penrose pseudoinverse). For $g=\vecop(G)$, $G\in\bR^{n\times r}$, the
preconditioned vector $z=P_{\mathrm{c}}^{\,\dagger}g$ that CG requires is
\begin{equation}\label{eq:precapply}
  z=\vecop\!\big(K^{+}\,G\,C^{-1}\big),
\end{equation}
where $C^{-1}\in\bR^{r\times r}$ is computed once. When $K$ is nonsingular,
$K^{+}=K^{-1}$ and \eqref{eq:precapply} solves $P_{\mathrm{c}}z=g$ exactly; when $K$
is singular, \eqref{eq:precapply} returns the solution in $\range(I_r\otimes K)$,
which is all CG ever needs.
\end{proposition}
\begin{proof}
By \eqref{eq:vec}, $P_{\mathrm{c}}\vecop(X)=(C\otimes K)\vecop(X)=\vecop(KXC^{T})=\vecop(KXC)$
since $C=C^{T}$. Hence solving $P_{\mathrm{c}}z=g$ means $KXC=G$ with $z=\vecop(X)$.
If $K$ is nonsingular, $X=K^{-1}GC^{-1}$. In general, the Moore--Penrose pseudoinverse
of $C\otimes K$ is $C^{-1}\otimes K^{+}$ (using $(A\otimes B)^{\dagger}=A^{\dagger}\otimes B^{\dagger}$
and $C^{\dagger}=C^{-1}$), and $(C^{-1}\otimes K^{+})\vecop(G)=\vecop(K^{+}GC^{-1})$
by \eqref{eq:vec}. Because $g=\vecop(G)$ produced inside CG always lies in
$\range(I_r\otimes K)$ (residuals of \eqref{eq:system} stay in $\range H$), the
pseudoinverse application returns the unique preimage in $\range(I_r\otimes K)$,
consistent with the subspace on which CG runs. All operations in \eqref{eq:precapply}
($K^{+}$ via the $n\times n$ eigendecomposition, the $n\times r$ by $r\times r$ and
$n\times n$ by $n\times r$ products) involve only $K$, $C$, $G$, never $M$ or $N$.
\end{proof}

\section{Algorithms}

\subsection{Direct symmetric method (baseline)}

We first record an explicit assembly of \eqref{eq:system} that already avoids $O(M)$
and $O(N)$ work, then solve by a dense symmetric factorization. By
Proposition~\ref{prop:matvec}'s derivation, the rows of $S^{T}(Z\otimes K)$ are the
$q$ Kronecker rows $\tZ[\ell,:]\otimes\tK[\ell,:]$ (this is the row of $Z\otimes K$
at vec-index $i_\ell+n(j_\ell-1)$), so
\[
  H=\Phi^{T}\Phi+\lambda(I_r\otimes K),\qquad
  \Phi\in\bR^{q\times nr},\quad \Phi[\ell,:]=\tZ[\ell,:]\otimes\tK[\ell,:].
\]

\begin{algorithm}[H]
\caption{Direct symmetric solve of the mode-$k$ RKHS subproblem}
\begin{algorithmic}[1]
\Require kernel $K\in\bR^{n\times n}$ (sym.\ PSD), observed rows $\tZ\in\bR^{q\times r}$,
row indices $\{i_\ell\}_{\ell=1}^q$, observed values $\{t_\ell\}$, $\lambda>0$
\Ensure $W\in\bR^{n\times r}$ solving \eqref{eq:system}
\State $\tK[\ell,:]\gets K[i_\ell,:]$ for $\ell=1,\dots,q$ \Comment{observed kernel rows, $O(qn)$}
\State Form $\Phi\in\bR^{q\times nr}$ with $\Phi[\ell,:]\gets \tZ[\ell,:]\otimes\tK[\ell,:]$ \Comment{$O(qnr)$}
\State $H\gets \Phi^{T}\Phi+\lambda(I_r\otimes K)$ \Comment{$O(q\,n^2r^2)$}
\State $B\gets \sum_{\ell=1}^q t_\ell\,e_{i_\ell}\,\tZ[\ell,:]$;\quad $g\gets\vecop(KB)$ \Comment{Prop.~\ref{prop:rhs}, $O(qr+n^2r)$}
\If{$K$ nonsingular} \State solve $H\,\vecop(W)=g$ by Cholesky \Comment{$O(n^3r^3)$}
\Else \State solve in $\range(I_r\otimes K)$ by a symmetric (pseudo)inverse / LDL$^{T}$ \Comment{$O(n^3r^3)$}
\EndIf
\State \Return $W$
\end{algorithmic}
\end{algorithm}

\paragraph{Correctness and cost.}
Steps~2--3 reproduce $H$ exactly by the row identity above; step~4 is
Proposition~\ref{prop:rhs}. The solve is exact (Lemma~\ref{lem:psd}) or returns the
minimum-norm solution (Lemma~\ref{lem:consistent}). No step references the $M$
unobserved columns or the $N$-dimensional vectorization, so the cost is
$O\!\big(q\,n^2r^2+n^3r^3\big)$ time and $O(n^2r^2)$ memory for $H$. The cubic term
$O(n^3r^3)$ is the stated baseline; assembling $H$ alone is the additional expense the
problem flags.

\subsection{Preconditioned conjugate gradients}

\begin{algorithm}[H]
\caption{PCG for the mode-$k$ RKHS subproblem (matrix-free)}
\begin{algorithmic}[1]
\Require kernel $K\in\bR^{n\times n}$ (sym.\ PSD), observed rows $\tZ\in\bR^{q\times r}$,
row indices $\{i_\ell\}_{\ell=1}^q$, observed values $\{t_\ell\}$, $\lambda>0$,
tolerance $\epsilon$, max iterations $L$
\Ensure $W\in\bR^{n\times r}$ approximately solving \eqref{eq:system}
\Statex \textbf{Setup (once):}
\State $K=V\Sigma V^{T}$ (symmetric eigendecomposition); $K^{+}=V\diag(\sigma^{+})V^{T}$ \Comment{$O(n^3)$}
\State $C\gets \tZ^{T}\tZ+\lambda I_r$;\quad factor $C=L_cL_c^{T}$ (Cholesky) \Comment{$O(qr^2+r^3)$}
\State $B\gets \sum_{\ell=1}^q t_\ell\,e_{i_\ell}\,\tZ[\ell,:]$;\quad $g\gets\vecop(KB)$ \Comment{$O(qr+n^2r)$}
\Statex \textbf{Matrix-free operators} (Props.~\ref{prop:matvec},\ref{prop:precsolve}):
\Statex \quad $\mathrm{Hmv}(W)$: $G\gets KW$; $m_\ell\gets G[i_\ell,:]\cdot\tZ[\ell,:]$;
$P\gets\sum_\ell m_\ell e_{i_\ell}\tZ[\ell,:]$; \textbf{return} $KP+\lambda KW$.
\Statex \quad $\mathrm{Prec}(G)$: \textbf{return} $K^{+}\,G\,C^{-1}$ (apply $C^{-1}$ via $L_c$).
\Statex \textbf{PCG iteration} (all variables are $n\times r$ matrices, $\langle A,B\rangle=\sum_{ij}A_{ij}B_{ij}$):
\State $W\gets 0$;\ \ $R\gets \mathrm{reshape}(g)$;\ \ $Y\gets\mathrm{Prec}(R)$;\ \ $D\gets Y$;\ \ $\rho\gets\langle R,Y\rangle$
\For{$t=1$ \textbf{to} $L$}
  \State $Q\gets \mathrm{Hmv}(D)$;\quad $\alpha\gets \rho/\langle D,Q\rangle$
  \State $W\gets W+\alpha D$;\quad $R\gets R-\alpha Q$
  \If{$\|R\|_F\le \epsilon\,\|g\|$} \textbf{break} \EndIf
  \State $Y\gets \mathrm{Prec}(R)$;\quad $\rho'\gets\langle R,Y\rangle$;\quad $\beta\gets \rho'/\rho$
  \State $D\gets Y+\beta D$;\quad $\rho\gets\rho'$
\EndFor
\State \Return $W$
\end{algorithmic}
\end{algorithm}

\paragraph{Correctness.}
$\mathrm{Hmv}$ computes $H\vecop(W)$ exactly (Proposition~\ref{prop:matvec}) and
$\mathrm{Prec}$ applies $P_{\mathrm{c}}^{\dagger}$ exactly
(Proposition~\ref{prop:precsolve}). The body is textbook preconditioned CG with the
SPD preconditioner $P_{\mathrm{c}}$ restricted to $\range H$, on which both $H$ and
$P_{\mathrm{c}}$ are positive definite (Lemmas~\ref{lem:psd},\ref{lem:consistent},
\eqref{eq:precapply}). Hence the iteration converges to the unique solution in
$\range(I_r\otimes K)$, the minimum-norm solution of \eqref{eq:system}
(Lemma~\ref{lem:consistent}); when $K$ is nonsingular this is the unique exact
solution. Convergence in exact arithmetic occurs in at most
$\dim\range H=r\cdot\mathrm{rank}(K)\le nr$ steps, and the energy-norm error contracts
as $2\big(\frac{\sqrt\kappa-1}{\sqrt\kappa+1}\big)^{t}$ with
$\kappa=\kappa(P_{\mathrm{c}}^{\dagger}H|_{\range H})$, the preconditioned condition
number on the active subspace.

\paragraph{Termination.}
The {\bf for} loop runs at most $L$ iterations and exits earlier once the relative
residual falls below $\epsilon$; both the eigendecomposition and Cholesky in the setup
terminate (finite direct factorizations), so the algorithm always halts.

\paragraph{Complexity.}
\emph{Setup:} $O(n^3)$ (eigendecomposition of $K$) $+\,O(qr^2+r^3)$ (form/factor $C$)
$+\,O(qr+n^2r)$ (RHS) $=O(n^3+qr^2)$ (using $r<q$).
\emph{Per iteration:} $\mathrm{Hmv}$ costs $O(n^2r)$ for $KW$ and $KP$, plus $O(qr)$
for the gather/scatter through the observations, i.e.\ $O(n^2r+qr)$. $\mathrm{Prec}$
costs $O(nr\cdot\min(n,r))$ for $K^{+}G$ via $V$ and $O(nr^2)$ for $GC^{-1}$, i.e.\
$O(n^2r+nr^2)$. Inner products and {\sc axpy}s cost $O(nr)$. Thus one iteration is
$O(n^2r+nr^2+qr)$.
\emph{Total} for $L$ iterations:
\[
  O\!\big(n^3+qr^2\big)\;+\;L\cdot O\!\big(n^2r+nr^2+qr\big).
\]
Memory is $O(n^2+qr)$: the $n\times n$ factors of $K$, the $r\times r$ factor of $C$,
the $q\times r$ matrix $\tZ$, the index list, and a constant number of $n\times r$
work matrices. No object of size $\Theta(M)$, $\Theta(N)$, or $\Theta(nr\times nr)$ is
ever formed.

\paragraph{Comparison.}
The direct solver costs $O(n^3r^3)$ (Cholesky on the dense $nr\times nr$ matrix) and
$O(n^2r^2)$ memory. PCG replaces the cubic-in-$nr$ factorization by $L$ products of
cost $O(n^2r+nr^2+qr)$ after an $O(n^3+qr^2)$ setup. Because the preconditioner
\eqref{eq:precond} clusters the spectrum (it captures the exact regularizer and the
Kronecker backbone of the Gram term), $L$ is small and essentially independent of
$n,r,q$ for fixed problem geometry, so the iterative method is asymptotically cheaper
by roughly a factor $\Theta(nr)$ in time and $\Theta(nr)$ in memory.

\begin{remark}[Why only $\tZ,\tK$ are needed]
Every appearance of $Z$ in \eqref{eq:system} is sandwiched by $S^{T}$ (in the model
evaluation) or produces an $n\times r$ output via $R_yZ$ (in the adjoint), and in both
cases only the observed rows $\tZ$ survive (Lemmas~\ref{lem:model},\ref{lem:adjoint});
the kernel enters either as the small $n\times n$ matrix $K$ or, in the direct
assembly, through its observed rows $\tK$. This is exactly why the full $M\times r$
matrix $Z$ and the $N$-vector $\vecop(T)$ never need to be touched.
\end{remark}

\end{document}
